\documentclass[11pt,a4paper]{ctexart}
\usepackage[margin=2.4cm]{geometry}
\usepackage{amsmath,amssymb}
\usepackage{fancyhdr}
\usepackage{enumitem}
\linespread{1.18}
\setlength{\headheight}{15pt}
\setlength{\parskip}{0pt}

\pagestyle{fancy}
\fancyhf{}
\fancyhead[L]{\small Team 2212245}
\fancyhead[R]{\small \thepage}
\renewcommand{\headrulewidth}{0.4pt}
\fancypagestyle{plain}{\fancyhf{}\renewcommand{\headrulewidth}{0pt}}

\ctexset{
  section/format = \Large\bfseries\raggedright,
  subsection/format = \large\bfseries\raggedright,
  subsubsection/format = \normalsize\bfseries\raggedright
}
\setcounter{secnumdepth}{3}

\begin{document}

%====================== 摘要页 ======================
\thispagestyle{empty}
\begin{center}
{\small 2022 MCM/ICM\quad 选题 A\quad 队伍控制号 2212245\quad 摘要页（Summary Sheet）}\\[18pt]
{\LARGE\bfseries 如何完成一场自行车计时赛}\\[8pt]
{\large （How to Ride a Cycling Time Trial）}\\[10pt]
\end{center}

\begin{center}
{\large\bfseries 摘\quad 要}
\end{center}

\vspace{6pt}
「作为一名数学家，你习惯于独立地解决问题，我正是以这样的方式来对待自行车运动的。」
\begin{flushright}
——安娜·基森霍夫（Anna Kiesenhofer），2020 年东京奥运会女子公路赛冠军
\end{flushright}

正如拥有数学博士学位的安娜·基森霍夫所说，一个恰当的数学模型能够帮助自行车运动员更好、更高效地规划自己的比赛。而这恰恰就是本文所完成的工作。

我们首先对自行车运动员的能量供给系统进行分析。依据基础的生物学知识，我们把总输出功率的贡献划分为三种不同的能量供给系统：ATP-CP 系统、糖解系统（糖酵解系统）与有氧系统。我们以常微分方程为基础建立模型，用以描述骑行过程中人体内各种物质含量的变化。

在能量供给模型的基础上，我们进而定义了计时赛专精选手与冲刺型选手的功率曲线（功率档案），并分别考虑了不同性别的情形。将我们得到的理论功率曲线与真实世界的功率曲线相比较，结果表明我们所建立的能量模型的可靠性。值得注意的是，本文中的功率曲线只是作为一项中间结果给出；我们使用更为本质、也更为精确的能量模型来展开后续讨论。

随后，为评估维持给定速度所需的功率，我们分析了由自行车与运动员构成的系统的运动过程。我们把各类阻力、丘陵地形以及急转弯都纳入模型之中，并由此得到运动方程，用以确定加速度与推进力之间的关系。

在能量模型与运动模型建立之后，我们讨论了各类特征地形上的最优策略。我们发现，在大多数情形下最优策略是在整个比赛过程中保持一个固定速度，即阈值速度；例外情形出现在丘陵赛道或急弯道路上——上坡骑行或通过急弯时通常需要减速。据此，我们提出了一套可行的策略，它是各个具体地形上最优策略的组合。

为在真实世界条件下检验我们的模型，我们首先重建赛道，并通过处理官方路线图获取必要数据。随后，我们模拟了四类不同运动员在三条赛道上的实际表现。模拟结果显示，2021 年在比利时举行的国际自行车联盟（UCI）世界锦标赛赛道上的最优成绩为 53.02 分钟，与该赛道冠军 57.78 分钟的成绩非常接近。我们评估了策略偏离所带来的敏感性，发现 30\% 的偏离只会造成 3.11\% 的时间差异。至于天气的敏感性，我们着重考察了风向与风力。分析表明，风会极大地影响所采用的策略与预期表现。

最后，我们对模型进行修正，把团体计时赛纳入其中，在团体计时赛中考虑跟骑效应以降低整体空气阻力。我们提出了团体计时赛的策略：运动员轮转领骑，并在冲刺阶段之前及冲刺阶段之中调整自身状态。与个人情形相比，团队策略可以把速度提高 14.43\%。

\vspace{6pt}
\noindent\textbf{关键词：}自行车运动员；功率曲线；计时赛；常微分方程

\newpage
\pagestyle{fancy}

%%====================== 正文 ======================
\section{引言}

\subsection{背景}

在卡尔·冯·德莱斯（Karl von Drais）造出世界上第一辆自行车之后，骑行很快便成为一项广受欢迎的体育活动。自 1894 年首届现代奥林匹克运动会举办以来，自行车项目便出现在此后每一届奥运会之中[1]。如今，美国有超过 4400 万人参与公路自行车赛事，而且这一数字仍在持续增长[2]。2020 年东京奥运会的公路自行车比赛即以富士赛道（Fuji Speedway）为赛场[3]，该赛道依托赛车场的既有路面与设施布置，形成了弯道密集、坡度变化频繁的路线，这使得计时赛中的速度控制问题显得格外突出。

这条赛道在形态上具有典型的「快慢交替」特征：一方面有长度可观、可以长时间保持高速的主直道；另一方面又有若干半径很小的连续弯角，迫使运动员反复制动、反复再加速。就本文的问题而言，这种结构意味着最优策略不可能是单一的恒速骑行，而必须是「直道恒速、弯前减速」的分段组合，这也正是后文策略模型所要处理的核心情形。

公路自行车比赛有多种类型，其中我们重点关注计时赛项目。计时赛与普通的大组赛在本质上是不同的：大组赛中，运动员的主要对手是身边的其他人，战术的核心在于位置、卡位与终点前的冲刺；而在计时赛中，运动员几乎是在与时间本身较量，胜负取决于能否在给定的赛道上合理地分配自己的体力。正是这一特征，使计时赛成为最适合用数学模型加以刻画的自行车项目，也正是本文选择它的原因。计时赛主要有两种形式：

在个人计时赛（individual time trial，ITT）中，运动员逐个出发完成一段固定的赛道，他们被禁止相互配合或彼此近距离骑行。每位运动员完成赛道的时间都会被记录，用时最少者赢得比赛。

在团体计时赛（team time trial，TTT）中，由运动员组成的车队逐个与时间赛跑。与不允许配合的个人计时赛不同，在团体计时赛中，同一车队的运动员可以骑在其他队员的尾流之中，以降低空气阻力。团体计时赛的获胜队伍通常通过比较各队第 4 名运动员的完赛时间来确定[4]。这一判定方式意味着团队必须作为一个整体前进，任何一名队员的掉队都会直接损害全队的成绩，因而队形、领骑时长与轮转节奏都需要精心设计。

综上，无论是个人计时赛还是团体计时赛，其核心问题都是同一个：在给定赛道、给定天气与给定运动员能力的条件下，如何把有限的体力分配到赛道的每一处，使完赛时间最短。本文正是围绕这一问题展开建模、求解与检验。

\subsection{问题重述}

本文旨在建立一个恰当的模型，为自行车运动员与车队体育主管（Directeur Sportif）提供最优比赛策略以及可能的成绩预测。

\begin{itemize}
\item 建立适用于任意类型自行车运动员的模型，以确定其在整个赛道上的功率分配。换言之，给定赛道上的某一位置，模型应当给出建议施加的最优功率。该模型应当把运动员的功率曲线、天气状况与地形条件都纳入考虑。这一要求意味着模型的输出不是一个单一的数字，而是沿赛道位置变化的一条功率曲线，且该曲线必须与运动员的生理能力相容。
\item 为两类运动员、并区分不同性别，定义其功率曲线。我们的功率曲线应当包含计时赛专精选手的功率曲线。功率曲线是连接生理能力与比赛策略的桥梁：它把「这位运动员有多强」转换成「他能够在多长时间内输出多大功率」。
\item 把我们的模型应用于各种计时赛赛道，其中包括 2021 年的两条真实赛道，以及一条自行设计的、包含足够多急弯与爬坡的赛道。真实赛道用于检验模型与现实的一致性，自行设计的赛道则用于在参数可控的条件下揭示模型的行为。
\item 针对可能偏离最优功率分配的情形，以及天气与环境的小幅变化，进行敏感性分析。这一步的目的在于判断我们的策略是否具有实用价值：若运动员执行策略时存在偏差，成绩会损失多少；若风向与风力发生变化，最优策略又应当如何调整。
\item 把我们的模型推广到团体计时赛，在该情形下运动员可以借助队友的帮助以获得更高的功率与更快的恢复。团体计时赛引入了一个新的机制——跟骑，它改变了阻力结构，也改变了体力在队员之间的分配方式。
\end{itemize}

\section{假设与符号说明}

\subsection{基本假设与理由}

\begin{itemize}
\item \textbf{把运动员与其自行车视为质点。}一场典型的奥运会计时赛赛道，女子约为 20 km 长，男子约为 40 km 长。在如此宏大的尺度上，运动员自身的尺寸是可以忽略的。这一假设使我们能够用位置与速度这两个量完整地描述运动员的状态，而不必引入姿态与转动惯量等难以标定的量；其代价是无法刻画起步、摇车等细节，但这些细节对总时间的影响远小于赛道尺度上的功率分配。
\item \textbf{假设赛道中的急转弯路段相对平坦，且运动员在急转弯中保持匀速运动。}出于安全考虑，赛事组织者通常会为急弯选择相对平坦的区域。因此在弯道内部，我们可以把运动近似为半径固定的匀速圆周运动，从而使过弯速度上限可以由向心力与摩擦力的平衡直接给出。
\item \textbf{同一性别的运动员具有相同的体重，即该性别的平均体重。}尽管运动员体重是影响问题结果的一个因素，但它并不是我们需要探究的主要问题，因此可以将其忽略。进一步地，我们所有的功率数值都以单位体重的形式给出，这样模型结果可以直接按体重换算到具体个体，而不必为每一位运动员重新标定参数。
\item \textbf{人体能量供给的生物学过程以恒定速率进行。}事实上，人体能量供给的生物学过程非常复杂，其速率受到酶浓度等诸多因素的影响，难以给出数学表达式。通过这一假设，能量供给问题的分析在短时间内变得可行。该假设的合理性在于：我们关心的是数分钟到一小时尺度上的功率上限，而非毫秒尺度上的反应动力学，因而把各条通路的速率取为其在当前浓度下的平均值，足以刻画宏观的供能能力。
\end{itemize}

\subsection{符号说明}

我们把全文使用的主要符号按用途分为三组：能量供给相关符号、运动学相关符号与阻力相关符号。

能量供给相关的符号为：
\begin{itemize}
\item $v_{ATP,i}\ (i=1,2,3)$：各能量供给系统的 ATP 生成速率。
\item $\tilde{v}_{ATP,i}\ (i=1,2,3)$：各能量供给系统的最大 ATP 生成速率。
\item $v_{ATP}$：总体的 ATP 生成速率。
\item $\mu$：单位 ATP 所含的能量。
\item $P$：运动员的总输出功率。
\item $P_i\ (i=1,2,3)$：各部分所对应的输出功率。
\item $V_x\ (x=CP,O_2,LA)$：相应物质的含量。
\item $\alpha_i\ (i=1,2,3)$：三种能量供给方式下的转化效率。
\item $M_x\ (x=CP,O_2,LA)$：相应物质的储存上限。
\end{itemize}

运动学与阻力相关的符号为：
\begin{itemize}
\item $f_a$：空气阻力。
\item $f_r$：滚动阻力。
\item $f_s$：坡度阻力。
\item $m^{*}$：自行车—运动员系统的有效质量。
\item $m$：自行车—运动员系统的质量之和。
\item $p$：空气阻力系数。
\item $\eta$：传动系统效率。
\item $F$：推进力。
\item $v_r$：相对速度。
\end{itemize}

此外，后文在讨论赛道与天气时还会用到弯道半径 $r$、路面摩擦系数 $k$、坡面倾角 $\alpha$、车身倾角 $\beta$ 以及风速分量 $v_x$、$v_y$ 等符号，它们均在首次出现处给出定义。

\section{自行车运动员的能量供给模型：三种能量供给系统的混合}

每位运动员的功率曲线都是独一无二的，这是由他们能量供给系统各自的特性所决定的。举例来说，冲刺型运动员具有更强的爆发力，其原因在于 ATP-CP 能量供给系统的性能更佳，而该系统正是短时间、极高强度运动的主要能量供给系统。反之，计时赛专精选手在有氧系统上的能力更为突出，这使他们在长达数十分钟的稳定输出中不易衰竭。由此可见，只要能够把三种供能系统各自的容量与速率刻画清楚，运动员的能力结构就可以被定量地描述，而功率曲线不过是这一描述的外在表现。

众所周知，ATP 是人体直接的能量来源。肌肉收缩所需的能量并不直接来自糖或脂肪，而是来自 ATP 水解为 ADP 与无机磷酸时释放的化学能。因此，无论哪一种供能通路，其最终作用都是把 ADP 重新磷酸化为 ATP。ATP 的生成速度与输出功率通过一个常数相互联系，即单位 ATP 所含的能量 $\mu$。于是输出功率 $P$ 与 ATP 生成速率 $v_{ATP}$ 满足如下关系式：
\begin{equation}
P=\mu v_{ATP}.
\end{equation}
这一关系式是整个模型的枢纽：它把生物化学层面的「物质周转速率」与力学层面的「功率」联系起来，使我们能够用描述化学物质含量变化的微分方程直接推算运动员的输出能力。

一般而言，细胞中仅储存有少量的 ATP，在剧烈运动中只能维持 1--3 s。其后继的能量供给主要依赖 ATP 的再合成。这就解释了为什么任何一次爆发都不可能无限持续：可用的 ATP 储备在数秒之内便会被耗尽，此后输出功率的高低完全取决于再合成速率。

人体内共有三种能够生成 ATP 以提供能量的能量供给系统，它们分别是 ATP-CP 系统、糖解系统与有氧系统[5]。总体输出功率被视为各能量供给系统输出功率之和：
\begin{equation}
P=P_1+P_2+P_3.
\end{equation}
需要强调的是，这一叠加并不是说三种系统各自独立地「分工」，而是说在任一瞬间，总功率由三者按当前生理状态所决定的比例共同承担；随着时间推移，比例会自发地由快速通路向慢速通路转移。

下面我们逐一讨论各能量供给系统的 ATP 生成过程。

\subsection{ATP-CP 系统}

细胞内磷酸肌酸（creatine phosphate，CP）的分解能够提供能量与无机磷酸，从而使 ADP 得以再合成 ATP。这一过程的机理可概括为：CP 在肌酸激酶的催化下水解，释放出的高能磷酸基团转移至 ADP，重新生成 ATP，同时产生肌酸；由于该反应无需氧气参与、且不经过一系列复杂的酶促中间步骤，因而可以在极短的时间内以极高的速率进行，这正是 ATP-CP 系统能够支撑爆发性用力（如起跑、冲刺与陡坡强攻）的生物学基础。与之相对应，该系统的总容量十分有限，一旦 CP 储备被消耗殆尽，功率便必须由其他通路接管。

该系统的示意可以概括为一条单向的反应链：CP 在酶的催化下分解，释放出高能磷酸基团与无机磷酸，后者与被肌肉做功「用掉」的 ADP 结合，重新生成 ATP；与此同时，分解产物肌酸在细胞中积累。与之相反，恢复过程依赖有氧代谢所提供的能量，把肌酸重新磷酸化为 CP，从而为下一次爆发储备原料。图中标注的三类物质（CP、ADP/ATP 与肌酸）此消彼长，正是 (3) 式中三个方程所描述的对象：CP 单调下降、ATP 以正比于 CP 分解速率的速率生成、初始时刻 CP 处于满储存。

\subsubsection{ATP-CP 系统的 ATP 生成}

对于 CP 的消耗过程，我们可以写出
\begin{equation}
\begin{cases}
\mathrm{d}V_{CP}(t)=-v_{CP}\,\mathrm{d}t,\\[0.3em]
v_{ATP,1}=\alpha_1 v_{CP},\\[0.3em]
V_{CP}(t=0)=M_{CP},
\end{cases}
\end{equation}
其中 $V_{CP}$ 为人体内的 CP 浓度，$M_{CP}$ 为人体内 CP 的最大储存量，$\alpha_1$ 为分解单位 CP 所生成的 ATP 数量，$v_{CP}$ 为 CP 的分解速率。第一式描述 CP 含量随时间的线性衰减，第二式说明 ATP 的生成速率正比于 CP 的分解速率，第三式则给出初始时刻 CP 处于满储存状态。把第一式积分，可得 $V_{CP}(t)=M_{CP}-v_{CP}t$，即 CP 的剩余量随骑行时间线性减少；当 $V_{CP}$ 降为零时，该系统再也不能提供功率。

求解上述方程组，我们得到
\begin{equation}
\alpha_1 V_{CP}(t)=v_{ATP,1}\tau,
\end{equation}
其中 $\tau$ 为给定 ATP 生成速率所能持续的时间。该式表明，在 CP 分解速率恒定的假定下，可生成的 ATP 总量与持续时间成正比，而单位时间内可供输出的 ATP 量则随 CP 储存的消耗而线性下降；等价地，若要求在时间 $\tau$ 内维持某个 ATP 生成速率，则所需的 CP 存量正比于 $\tau$，这意味着持续时间越长，可用的平均功率越低。

由于实际生物结构的限制，ATP 的生成速度必然受到生理学上的约束。设 ATP-CP 系统的最大 ATP 生成速度为 $\tilde{v}_{ATP,1}$，利用输出功率与 ATP 生成速率之间的关系式，我们得到
\begin{equation}
P_1=\mu v_{ATP,1}=\mu\min\left\{\frac{\alpha_1 V_{CP}}{\tau},\ \tilde{v}_{ATP,1}\right\}.
\end{equation}
式中取最小值这一步具有明确的生理含义：一方面，输出功率不能超过酶促反应所能达到的速率上限；另一方面，输出功率也不能超过在给定持续时间内可供消耗的 CP 存量所允许的水平。前者是「瞬时能力」的限制，后者是「续航能力」的限制，二者共同决定了该通路在任一时刻可贡献的功率。

ATP-CP 过程在有氧与无氧条件下均可进行，持续时间约为 6--8 s。其输出功率非常高，因此它是短时间、高强度运动的主要能量供给系统。

\subsubsection{CP 的恢复}

在有氧呼吸过程中，人体内的 CP 浓度得以恢复。若把 CP 的恢复速率定义为 $v_{r1}$，则恢复过程的方程与初始条件为
\begin{equation}
\begin{cases}
\mathrm{d}V_{CP}(t)=v_{r1}\,\mathrm{d}t,\\[0.3em]
V_{CP}(t=0)=0.
\end{cases}
\end{equation}

当 $V_{CP}$ 达到其最大值 $M_{CP}$ 时，恢复过程结束。在完全耗尽之后，CP 的恢复时间约为 3--5 min。这一时间尺度在策略设计中具有直接意义：它决定了运动员能够在多短的间隔内重复一次高强度的 ATP-CP 供能，也决定了运动员在赛道上一次全力强攻之后需要多长的相对轻松路段才能重新获得爆发能力。正因如此，合理的策略不会把这一通路用在赛道的开头，而是把它保留到最后数百米。

\subsection{糖解—有氧系统}

糖解系统与有氧系统彼此紧密协作，因此我们在本小节中把它们合并考虑。

在无氧条件下，糖原或葡萄糖经糖解作用被降解，产生乳酸与 ATP，其中大部分能量储存于乳酸之中。该过程同时也是葡萄糖氧化代谢的准备途径。在有氧条件下，乳酸继续通过氧化代谢分解，生成大量 ATP 以提供能量。整体过程可概括为：糖原或葡萄糖在细胞质中经一系列酶促反应分解为丙酮酸，并进一步还原为乳酸，同时净生成少量 ATP；乳酸随后被转运至氧化代谢活跃的肌纤维或其他组织，在有氧条件下进入三羧酸循环与电子传递链，被彻底氧化为二氧化碳与水，并释放出大量能量用于 ATP 的再合成。上述两个环节首尾相接，共同构成一条完整的供能通路：前者提供快速但「不彻底」的能量释放，后者提供缓慢但「彻底」的能量释放，二者通过乳酸这一中间产物相互衔接。

这一通路的示意可以概括为上下游两段：上游是糖原或葡萄糖经糖解降解为乳酸、同时净生成少量 ATP 的过程，其速率 $v_{ge}$ 决定了糖解系统的功率；下游是乳酸在有氧条件下被氧化、按化学计量关系消耗氧气并生成大量 ATP 的过程，其速率 $v_{de}$ 经系数 $k$ 换算为耗氧速率 $v_{out}$，进而决定有氧系统的功率。图中同时标出了氧的供给项 $v_{in}$，它代表肺与循环系统持续向肌肉输送氧气的能力，正是这一项使有氧系统能够在长时间内稳定工作。

\subsubsection{糖解—有氧系统的 ATP 生成}

与我们在 ATP-CP 系统中所做的处理类似，我们可以为糖解系统写出方程。
\begin{equation}
\begin{cases}
v_{out}=k v_{de},\\[0.3em]
\mathrm{d}V_{O_2}(t)=(v_{in}-v_{out})\,\mathrm{d}t,\\[0.3em]
\mathrm{d}V_{LA}(t)=(v_{ge}-v_{de})\,\mathrm{d}t,\\[0.3em]
v_{ATP,2}=\alpha_2 v_{ge},\\[0.3em]
v_{ATP,3}=\alpha_3 v_{out},\\[0.3em]
V_{O_2}(t=0)=M_{O_2},\\[0.3em]
V_{LA}(t=0)=0,
\end{cases}
\end{equation}
其中 $k$ 为分解单位乳酸所需的氧气量，$\alpha_2$、$\alpha_3$ 分别为生成单位乳酸或消耗单位氧气所生成的 ATP 数量，$V_{O_2}$ 为人体内的氧浓度，$M_{O_2}$ 为最大氧浓度，$V_{LA}$ 为乳酸浓度，$v_{in}$ 为肺向人体提供氧气的速率，$v_{out}$ 为有氧呼吸所消耗氧气的速率，$v_{ge}$ 为糖分解生成乳酸的速率，$v_{de}$ 为有氧条件下乳酸分解的速率。逐一来看：第一式把氧气消耗速率与乳酸分解速率通过化学计量系数 $k$ 联系起来；第二式与第三式分别是氧气与乳酸的守恒方程，说明其含量的变化率等于生成速率与消耗速率之差；第四式与第五式则表明糖解系统与有氧系统各自的 ATP 生成速率分别正比于糖解速率与耗氧速率；最后两式给出初始条件，即比赛开始时氧气储备处于上限、乳酸浓度为零。

求解上述方程组，我们得到
\begin{equation}
\begin{cases}
\alpha_2\left(M_{LA}-V_{LA}+\dfrac{V_{O_2}+v_{in}\tau}{k}\right)=v_{ATP,2}\tau,\\[0.9em]
\alpha_3\left(V_{O_2}+v_{in}\tau\right)=v_{ATP,3}\tau.
\end{cases}
\end{equation}

与 ATP-CP 系统的情形类似，ATP 的生成速度受到生理学上的限制，即 $v_{ATP,2}\leqslant\tilde{v}_{ATP,2}$ 且 $v_{ATP,3}\leqslant\tilde{v}_{ATP,3}$。再次利用输出功率与 ATP 生成速率的关系式，我们最终得到
\begin{equation}
P_2=\mu v_{ATP,2}=\mu\min\left\{\frac{\alpha_2\left(M_{LA}-V_{LA}+V_{O_2}/k\right)}{\tau}+\frac{v_{in}}{k},\ \tilde{v}_{ATP,2}\right\},
\end{equation}
\begin{equation}
P_3=\mu v_{ATP,3}=\mu\min\left\{\alpha_3\left(\frac{V_{O_2}}{\tau}+v_{in}\right),\ \tilde{v}_{ATP,3}\right\}.
\end{equation}

从这两个表达式中可以读出模型的两条重要性质：其一，糖解系统的可用功率取决于剩余的可容纳乳酸容量 $M_{LA}-V_{LA}$，当乳酸积累接近其阈值时，糖解系统所能贡献的功率迅速衰减；其二，有氧系统的功率由体内氧储备与肺的供氧速率共同决定，其中 $v_{in}$ 项对应稳态供氧能力，$V_{O_2}/\tau$ 项则对应储氧被一次性透支所能换取的额外功率，随着持续时间 $\tau$ 的延长，后一项迅速变得可以忽略，这与有氧供能适合长时间稳定输出的生物学事实一致。换言之，模型自动地重现了一条经验规律：冲刺看的是能够透支多少，计时赛看的是能够稳定供给多少。

\subsubsection{糖解子系统的恢复}

当乳酸浓度达到乳酸阈值 $M_{LA}$ 时，糖解过程结束。只有当糖解系统通过有氧呼吸分解乳酸而得到恢复之后，糖解过程才能继续进行。而当人体内的氧浓度不足时，有氧呼吸的强度同样会受到抑制。换言之，糖解系统与有氧系统在恢复环节上互为条件：糖解的恢复依赖有氧代谢清除乳酸，而有氧代谢的强度又受制于氧的供给。这一耦合关系解释了为何在高强度比赛中，运动员一旦使乳酸积累逼近阈值，就必须把功率降到有氧区间才能重新获得再次强攻的能力。

\subsection{整体能量供给：将三者合并}

真实世界中的人体能量供给是由三种能量供给系统同时完成的。某一系统是否占据主导地位，主要取决于运动的强度水平。强度较低时，有氧系统可以在不产生乳酸积累的前提下持续供能；强度提高后，糖解系统开始介入，乳酸随之积累；强度再高、时间再短，则由 ATP-CP 系统承担主要份额。三者的相对贡献随时间与强度的变化，构成了运动员功率曲线的全部内容。

整体供能的示意结构可以概括为「三进一出」：三条并联的通路各自把各自的底物转化为 ATP，其汇合点即总的 ATP 生成池，再由该池按 $P=\mu v_{ATP}$ 转换为输出功率；而在这个汇合点之前，还有一个限幅环节，它依据当前的 CP、乳酸与氧浓度算出这一时刻的功率上限，并把超限的请求折回为「不可实现」或以乳酸透支的方式记账。正是这一限幅环节，使得模型在短时间尺度上呈现出功率曲线的陡降段，在长时间尺度上则呈现出趋于平台的有氧极限。

首先要指出，在一场比赛的整个过程中，人体内的糖总量是固定的。然而在实际比赛中，运动员一般并不会出现糖的匮乏。真正限制运动员能量供给的是肌肉疲劳（可用乳酸与 CP 的浓度加以量化）以及其他问题。因此，糖总量这一限制在我们的分析中实际上并不会起作用。这一判断使模型得以简化：我们无需追踪糖原储量这一难以标定的状态量，只需追踪 CP、乳酸与氧三个状态量，便足以决定任一时刻的功率上限。

我们定义一个算法以确定占主导地位的能量供给过程：

\begin{enumerate}
\item 人体发出指令，要求输出一定的功率 $P$ 并持续一定的时间 $t$。
\item 根据持续时间 $t$ 以及人体内 CP、乳酸与氧的当前浓度，计算最大输出功率上限 $P_{limit}$。
\item 根据所需功率选择占主导地位的能量供给系统。当功率较低时，有氧系统即可满足功率需求；当有氧功率不足时，考虑加入糖解系统，乃至 ATP-CP 系统以供能。如果三种能量供给系统的功率总和仍然不足，则报告为不可实现。
\end{enumerate}

这一算法的实质是一个在每个时间步上执行的判据：先由当前状态算出这一时刻可以施加的最大功率，再由实际施加的功率反推各系统的消耗速率，进而更新状态量，如此循环推进。它把连续的生理过程离散为一系列可计算的状态转移，从而使整场比赛的仿真成为可能。

最后，考虑肾上腺素等激素的作用，上述模型还可以进一步修正。当三种能量供给系统共同工作仍无法满足功率输出需求时，允许通过糖解在有限范围内强制超出功率上限，其代价是产生更多的乳酸，即超出乳酸阈值的那一部分。引入超限常数 $e>1$，我们把总体输出功率的表达式修正为
\begin{equation}
P\leqslant e\left(P_1+P_2+P_3\right).
\end{equation}

与正常的糖解相比，超过乳酸阈值的糖解更加不彻底。因此，若引入描述分解不彻底的常数 $b$，则新产生的乳酸含量为
\begin{equation}
\mathrm{d}V_{LA}=b\left[P-\left(P_1+P_2+P_3\right)\right]\mathrm{d}t.
\end{equation}
该式刻画了「透支」的代价结构：超出供能能力的那部分功率并不会凭空产生，而是以更快的速率转化为乳酸的积累，从而在随后的赛程中以功率上限下降的形式被「偿还」。正是这一机制，使得运动员在终点前的短距离全力冲刺成为可能，同时也使得过早透支变得得不偿失。超限常数 $e$ 则刻画了在激素作用下人体所能承受的瞬时超额程度，它把最后冲刺这一现实中屡见不鲜的现象纳入了模型的描述范围。
\section{能量供给模型的应用：自行车运动员的功率曲线}

在建立了自行车运动员的能量供给模型之后，我们便可以深入讨论功率曲线。功率曲线之所以值得专门讨论，是因为它是教练与运动员在现实中真正能够测量、并且真正用来制定战术的对象；而我们的模型则提供了从生理机制出发解释并预测这条曲线的途径。

\subsection{基于能量供给模型求解功率曲线}

运动员的功率曲线描述的是其维持不同水平输出功率所能持续的时间长度，而这一曲线借助我们的能量供给模型很容易得到。

按照功率曲线的定义，总体功率曲线是三种能量供给系统各自功率曲线之和。求解时采用的初始条件为
\[
V_{CP}(t=0)=M_{CP},\qquad V_{O_2}(t=0)=M_{O_2},\qquad V_{LA}(t=0)=0.
\]
这三个初始条件分别对应三种物质在比赛开始时所处的状态：CP 与氧的储备处于上限，乳酸尚未开始积累。正是这一「净空」的初始状态，使运动员在起步阶段能够同时调动三条通路，从而在功率曲线上表现为最左端那个极高的尖峰。

在按照运动员的特征赋予上述初始值以及其他相关参数之后，把能量供给模型中三种能量供给系统的功率曲线表达式相加，即可得到该运动员的功率曲线。具体做法是：对每一个候选的持续时间 $\tau$，分别由三种供能系统的表达式算出该持续时间下所能维持的最大功率，再取三者之和，便得到功率曲线上与该 $\tau$ 对应的一个点；令 $\tau$ 由数秒变化到一小时以上，即可描出完整的曲线。由于三种系统所对应的表达式随 $\tau$ 的衰减速度差别极大，曲线在双对数坐标下呈现为几段斜率不同的近似直线，各段的斜率正对应着各系统的主导区间。

需要说明的是，由于实际生物学过程十分复杂，参数有意义的值是参数之间的相对值，而非绝对值；我们所得出的功率曲线与功率数值都已经过处理，以消除运动员体重的影响，也就是说，下文展示的功率数值乘以运动员的体重，才是实际的输出功率。这一处理使我们能够把不同性别、不同体重运动员的能力放在同一标尺下比较，也使模型结果可以方便地应用于具体个体：只要知道某位运动员的体重，就可以把单位体重的功率换算为他实际需要盯住的瓦特数。

\subsection{运动员的性别差异}

不同性别运动员的功率曲线特征也存在显著差异。

首先，我们引入功率曲线上最为重要的四个数值[6]：
\begin{itemize}
\item 冲刺能力（5 s 时的功率）；
\item 无氧能力（1 min 时的功率）；
\item 最大摄氧能力（5 min 时的功率）；
\item 功能阈值功率（1 h 时的功率）。
\end{itemize}

这四个点之所以被选作特征值，是因为它们分别落在三条供能通路各自的主导区间之内：5 s 对应 ATP-CP 系统，1 min 对应糖解系统，5 min 与 1 h 则对应有氧系统在短时透支与稳态输出两种状态下的能力。因此，这四个数值实际上就是对三种供能系统能力的四点采样。

FTP 即功能阈值功率（Functional Threshold Power），它估计的是运动员能够维持一小时的最高平均功率。在自行车运动中，FTP 是衡量体能水平的一项指标，它指示了你能够在长时间内维持多大的做功量[7]。

VO2Max 是位于「阈值」区间之上的第一个区间，它从 FTP 的 105\% 开始，到 120\% 结束[8]。VO2Max 方面的优势有助于运动员在短促而陡峭的爬坡中表现出色[9]。

我们将通过比较不同水平男女运动员的上述四个特征值，来展示不同性别运动员之间的差异。该分析基于从网络上检索到的数据[6]。比较的方式是：把上述四个特征值作为纵坐标、把运动员的竞技水平作为横坐标，在同一张图上同时呈现男女两类数据，其中实心圆点代表男性，虚线三角形代表女性；四种颜色分别对应四个重要的参考数值，横轴上的数字越小表示运动员的水平越高，纵轴则是消除体重影响之后的功率数值。

通过分析这些数据，不难得出如下结论：
\begin{itemize}
\item 在所有持续时间长度上，女性运动员的输出功率强度都低于男性运动员，而且随着运动员水平的提高，这一差异愈发明显；
\item 尽管同一水平女性运动员的能力普遍低于男性运动员，但其差异在短时爆发力方面更大，而在长时间耐力方面的差异相对较小。
\end{itemize}

上述结论也与男女两性的生理特征相一致。从模型的角度看，这两个结论分别对应两类不同的参数调整：短时爆发力的差异主要体现在 ATP-CP 系统的最大 ATP 生成速率与 CP 储存上限上，而长时间耐力的差异则主要体现在有氧系统的最大 ATP 生成速率与最大氧浓度上。这也提示我们，在把模型用于具体运动员时，性别并不是一个需要单独引入的修正项，它会自动地通过这两组参数的取值体现出来。

这一比较的图形呈现可以概括为四组「成对」的数据：同一颜色下的实心圆点与虚线三角形分别对应男女，两者的竖直间距即该能力上的性别差距；而随着横轴水平编号的减小（即水平提高），四组数据的整体高度都在上升，且男女之间的间距同时被拉大。图中最上方的两组对应短时能力（5 s 与 1 min），最下方的两组对应长时能力（5 min 与 1 h），四者的相对位置直观地显示出：短时能力的性别差距大于长时能力的性别差距，而所有能力的绝对水平都随训练水平提高而提高。

\subsection{计时赛专精选手的功率曲线}

我们基于 Python，使用 Numpy、Sympy 等库对我们的模型进行模拟、求解与分析。Numpy 用于向量化的数值积分与曲线绘制，Sympy 用于对方程组作符号求解，从而避免手工推导可能引入的错误。

计时赛专精选手的特征是耐力强、抗疲劳水平高，并且能够长时间保持接近阈值的功率。其重要参数的取值如下：对男性而言，CP 最大储存量 $M_{CP}=30$，肺供氧速率 $v_{in}=1.6$，最大氧浓度 $M_{O_2}=13000$，乳酸阈值 $M_{LA}=220$，ATP-CP 系统最大 ATP 生成速率 $\tilde{v}_{ATP,1}=20$，糖解系统最大 ATP 生成速率 $\tilde{v}_{ATP,2}=7.1$，有氧系统最大 ATP 生成速率 $\tilde{v}_{ATP,3}=5$；对女性而言，相应取值分别为 $M_{CP}=22$、$v_{in}=1.6$、$M_{O_2}=10500$、$M_{LA}=180$、$\tilde{v}_{ATP,1}=18$、$\tilde{v}_{ATP,2}=6.8$、$\tilde{v}_{ATP,3}=4.8$。这些参数的选取遵循两条原则：其一，它们应当使模型算出的四个特征功率落在真实计时赛选手的实测范围之内；其二，男女参数之间的比例关系应当与上一小节所观察到的性别差异方向一致。

把上述参数代入第 3.3 节所给出的算法，即可计算得到总体功率曲线，其中三条能量供给系统各自的功率曲线也一并给出；为使曲线起始阶段的差异显示得更为清晰，我们采用了双对数坐标。由计算结果可见，在极短的持续时间（数秒）内，曲线由 ATP-CP 系统主导并呈陡峭下降；在数十秒到数分钟的区间内，糖解系统成为主要贡献者；而在数分钟以上的长时间区间内，曲线逐渐趋于平缓，并由有氧系统单独支撑，其渐近水平即对应运动员的功能阈值功率。

从曲线的形态上可以读出三点：第一，总功率曲线在极左端（数秒）出现一个远高于其余部分的高峰，其数值由 ATP-CP 系统的最大生成速率封顶，此后迅速跌落；第二，在中段（数十秒至数分钟）曲线出现一个明显的「肩部」，这是糖解系统在乳酸尚未累积到阈值之前所能支撑的额外功率；第三，在右端（十分钟以上）曲线几乎水平，三条分曲线中的有氧项与总曲线重合，说明此时其余两条通路已不再贡献。男女两条总曲线的形状相同而整体高度不同，这与上一小节关于性别差异的结论一致。

计时赛专精选手的四个特征功率如下：以每千克体重所对应的瓦特数（$\mathrm{W\cdot kg^{-1}}$）计量，男性在 5 s、1 min、5 min 与 FTP（1 h）四个持续时间上的取值分别为 18.10、9.17、5.83 与 5.07；女性在相应四个持续时间上的取值分别为 16.00、8.17、5.47 与 4.58。从这组数字可以看出，男性相对女性的优势在 5 s 处约为 13\%，而在 1 h 处则收窄到约 10\%，这正是性别差异随持续时间延长而缩小的定量表现。

以下结果我们将采用与本小节类似的方式展示，后文不再重复说明。

\subsection{冲刺型选手的功率曲线}

冲刺型选手的特征是爆发力强，能够在短时间内产生很高的功率，因此他们擅长冲刺。与大组赛中的角色分工相对应，冲刺型选手的比赛方式通常是把体力保存到终点前的最后数百米，然后在极短的时间内倾泻而出。

对冲刺型选手，重要参数的取值为：男性 $M_{CP}=55$、$v_{in}=1.32$、$M_{O_2}=11000$、$M_{LA}=285$、$\tilde{v}_{ATP,1}=25$、$\tilde{v}_{ATP,2}=7.3$、$\tilde{v}_{ATP,3}=3.98$；女性 $M_{CP}=50$、$v_{in}=1.2$、$M_{O_2}=10000$、$M_{LA}=265$、$\tilde{v}_{ATP,1}=24$、$\tilde{v}_{ATP,2}=7.12$、$\tilde{v}_{ATP,3}=3.75$。与计时赛专精选手相比，冲刺型选手在 CP 储存量、乳酸阈值与 ATP-CP 系统的最大生成速率上都明显更高，而有氧系统的最大生成速率则更低，这一参数组合恰好刻画了「短时极强、长时较弱」的能力结构。

相应的四个特征功率为：男性在 5 s、1 min、5 min 与 FTP（1 h）上的取值分别为 22.28、9.65、5.11 与 4.07 $\mathrm{W\cdot kg^{-1}}$；女性分别为 20.87、9.00、4.80 与 3.84 $\mathrm{W\cdot kg^{-1}}$。把这些数值与计时赛专精选手的数值并置比较，可以看到一条清晰的分界线：在 5 s 处冲刺型选手高出约 4 $\mathrm{W\cdot kg^{-1}}$，而在 1 h 处反而低出约 1 $\mathrm{W\cdot kg^{-1}}$；两条曲线在数分钟附近相交，交点正是两类选手能力结构的分水岭。

与计时赛专精选手的曲线相比，冲刺型选手的曲线在左端更高、在右端更低，整体呈现出更陡的斜率：其在数秒处的峰值明显抬升，而在十分钟以上的平台则明显下沉。这一对比说明，所谓「选手类型」在模型中可以完全由参数组刻画，而不需要引入任何额外的分类机制——两类选手的差别，只是同一组方程在不同参数下的两种解。

我们的模型所给出的结果与真实世界的数据吻合良好，这证明了模型的可靠性。因此，我们可以由给定的运动员功率曲线反推三种能量供给系统的能力。借助我们的模型，我们还能够评估超出其功率曲线范围的爆发力以及恢复时间，而这正是选择策略的指导依据。换言之，功率曲线只是模型的一个「投影」，而模型本身还能回答功率曲线无法回答的问题：某位运动员在透支之后需要多久才能恢复到阈值输出，以及一次超出曲线范围的强攻能够换取多少时间上的收益。

\section{自行车运动员的运动模型}

在讨论赛道上的策略之前，我们需要建立一个描述运动员运动的模型。能量模型回答的是「运动员能够输出多少功率」，运动模型回答的则是「这些功率能够换取多少速度」，二者结合才能把生理能力翻译成完赛时间。

\subsection{骑行中的各类阻力}

运动员的运动受到多种阻力的影响，包括空气阻力、传动阻力、坡度阻力与滚动阻力。下面逐一分析各类阻力的影响。

\subsubsection{空气阻力}

对于在平坦道路上高速行驶的自行车而言，空气阻力是阻力的重要组成部分。事实上，当速度超过某个阈值之后，空气阻力会迅速成为阻力中的主导项，这也是计时赛中运动员普遍采用低风阻姿态与专用器材的原因。

运动员（以及自行车）所受的空气阻力可以分解为压差阻力与摩擦阻力。

当运动员与空气之间存在相对运动时，运动员与自行车会压缩前方的空气，这意味着运动员前方将出现一个压力相对较高的区域，而其后方由于空气的绕流逃逸将出现一个压力相对较低的区域。这两个部分之间的压力差便构成了压差阻力。与此同时，当运动员高速运动时，空气还会在自行车车身与人体的表面上产生摩擦阻力。在两分量之中，压差阻力通常占据主导地位，其大小与迎风面积和速度的平方密切相关；摩擦阻力则与物体表面的粗糙程度及空气的黏性有关，量级相对较小。

对空气阻力的实际讨论相当复杂。按照经典流体力学，空气阻力可以表示为如下公式[10]：
\begin{equation}
f_a=\frac{1}{2}C\rho S v_r^{2},
\end{equation}
其中 $f_a$ 为空气阻力，$C$ 为阻力系数，$\rho$ 为空气密度，$S$ 为运动员的迎风面积，$v_r$ 为运动员与空气之间的相对速度。考虑到对某一位特定的运动员而言，前三个参数在比赛过程中几乎保持不变，该公式可以改写为如下形式：
\begin{equation}
f_a=p v_r^{2},
\end{equation}
其中 $p=\frac{1}{2}C\rho S$ 为合并后的空气阻力系数。这一简化把空气阻力归结为一个可标定的单一系数与相对速度平方的乘积，从而使后续的运动方程得以解析处理；同时也说明，速度对空气阻力的影响是二次的，因此在高速段提高速度的代价远高于低速段。需要指出的是，这里的 $v_r$ 是相对速度而非地速，因此顺风与逆风会使同一功率对应的地速产生明显差异——这正是后文讨论天气影响时的出发点。

\subsubsection{传动阻力}

在骑行过程中，运动员所提供的能量有一部分因传动阻力而损耗，即传动系统各部件之间的摩擦。传动阻力所引起的能量损失可以用传动效率 $\eta$ 来表示。对于为比赛而设计的真实自行车，传动效率平均约为 0.95[11]。因此我们假设传动效率 $\eta=0.95$。这一取值意味着运动员所输出的功率中约有 5\% 无法到达路面，在长时间比赛中这项损失会累积成可观的时间代价。

\subsubsection{坡度阻力}

我们把自行车所受重力沿平行于坡面方向的分量视为阻力的一部分，即所谓的坡度阻力。设坡面的倾角为 $\alpha$，容易证明
\begin{equation}
f_s=mg\sin\alpha,\qquad \alpha\in\left[-\frac{\pi}{2},\frac{\pi}{2}\right],
\end{equation}
其中 $f_s$ 为坡度阻力，$\alpha$ 在上坡时为正、下坡时为负。该式的符号约定使上坡与下坡可以统一处理：上坡时坡度阻力与运动方向相反，需要由运动员额外做功加以克服；下坡时则为负，相当于重力在做功，运动员可以借此获得速度而不必输出功率。

\subsubsection{滚动阻力}

滚动阻力一般由车轮或地面的形变所引起，可以用如下方程表示：
\begin{equation}
f_r=C_{rr}N,
\end{equation}
其中 $f_r$ 为滚动阻力，$C_{rr}$ 为滚动阻力系数（rolling resistance coefficient，RRC），$N$ 为垂直于路面的法向力。普通自行车轮胎的滚动阻力系数在 0.007 至 0.014 之间变化[12]；而根据滑行实验，为比赛而设计的自行车其滚动阻力系数可低至 0.004[13]。

若运动员处于倾角为 $\alpha$ 的坡面上，则法向力可由 $N=mg\cos\alpha$ 给出，于是滚动阻力的公式可以写为
\begin{equation}
f_r=C_{rr}mg\cos\alpha.
\end{equation}
与空气阻力不同，滚动阻力与速度无关，只与法向力成正比，因此它构成了一项恒定的「基础税」：无论运动员骑得快还是慢，都必须为它付出功率。

\subsection{运动员的直线运动}

计入传动效率 $\eta$ 之后，运动员的功率方程为
\begin{equation}
\eta P=Fv,
\end{equation}
它给出了运动员的输出功率 $P$、运动员的当前速度 $v$ 与推进力 $F$ 之间的关系。该式表明，在功率固定的条件下，推进力与速度成反比：速度越高，可用于克服阻力的推进力越小，这正是高速段加速极为困难的原因。

此处的关键定理是动能定理，即动能的改变量等于外力所做的功[14]。动能是平动动能 $\frac{1}{2}mv^{2}$ 与每个车轮的转动动能 $\frac{1}{2}I\omega^{2}=\frac{1}{2}\left(I/R^{2}\right)v^{2}$ 之和，其中 $m$ 为运动员与自行车的质量，$R$ 为自行车车轮的外半径，$I$ 为车轮的转动惯量。由于车轮无滑动地滚动，其角速度满足 $\omega=v/R$，于是转动动能可以完全用平动速度表示。把两部分相加并注意共有两个车轮，我们得到
\begin{equation}
\frac{\mathrm{d}}{\mathrm{d}t}\left[\frac{1}{2}\left(m+\frac{2I}{R^{2}}\right)v^{2}\right]=\left(F-f_a-f_r-f_s\right)v.
\end{equation}
等式右端是各力所做功的功率：推进力做正功，三类阻力做负功，而传动损失已通过 $\eta$ 计入 $F$ 之中。

对上式作简单计算可得
\begin{equation}
m^{*}a=\left(m+\frac{2I}{R^{2}}\right)a=F-f_a-f_r-f_s,
\end{equation}
其中 $m^{*}=m+2I/R^{2}$ 为运动员与自行车的有效质量，$a$ 为运动员的加速度。有效质量的引入意味着车轮的转动惯量在加速时表现为额外的平动质量，因此启动阶段与加减速阶段的时间代价会比仅按整车质量估计的更高，这一细节在急弯前后的加减速过程中尤为明显。当 $a=0$ 时，该式退化为力的平衡条件，可用于求解给定功率下的稳态速度；而当 $a\neq0$ 时，它给出了一个一阶常微分方程，可用于推进整场比赛的数值仿真。

\subsection{运动员在弯道上的运动}

运动员沿弯道骑行时会向内倾斜以获得最大速度，这种方法被称为倾斜法。

运动员在通过弯道时使自行车保持在一条直线上，同时让自行车向弯道内侧倾斜。地面对自行车的支持力与自行车自身重力的合力，为自行车的圆周运动提供向心加速度。这一策略可以获得较高的过弯速度，同时也有利于观察前方的道路。

该过程的几何关系可以概括为：车身与竖直方向成倾角 $\beta$，支持力沿车身方向斜向上、重力竖直向下，二者的合力水平指向弯心。倾角越大，合力的水平分量相对于竖直分量越大，能够提供的向心加速度也越大；但倾角的上限由轮胎与路面的摩擦条件决定，一旦所需的水平力超过摩擦力所能提供的上限，倾斜便不再能换来更高的速度。图中以矩形示意自行车、以箭头示意支持力与重力的合成，其方向与圆周运动的向心方向一致。

按此方式，设自行车与竖直方向的倾角为 $\beta$，则支持力与重力的合力提供向心力，由此得到如下公式：
\begin{equation}
\frac{mv^{2}}{r}=\frac{mg}{\tan\beta}.
\end{equation}

考虑自行车轮胎与地面之间摩擦力的限制，该式应当写为
\begin{equation}
\frac{mv^{2}}{r}=\frac{mg}{\tan\beta}\leqslant kmg.
\end{equation}
其中的不等号来自轮胎与路面之间的静摩擦上限：一旦所需的向心力超过摩擦所能提供的最大值 $kmg$，轮胎便会侧滑。

因此，通过弯道所允许的最大速度为
\begin{equation}
v_m=\sqrt{grk},
\end{equation}
其中 $g$ 为重力加速度，$r$ 为弯道半径，$k$ 为地面摩擦系数，其值约为 0.3。当弯道的曲率半径较大时，这一限制变得相当宽松，在此类情形下运动员的运动近似等价于在直道上运动。这也正是后文在重建赛道时只把半径较小的弯道单独处理、而把半径较大的弯道视为直道的依据：以 $k=0.3$ 计，半径 60 m 的弯道所允许的最大速度约为 48 km/h，已经高于计时赛中的典型速度，因而对策略不再构成实质性约束；而半径更小的弯道则会强制压低速度，必须显式建模。
\section{个人计时赛的最优策略}

考虑运动员实际可能处于的状态，将包括以下三种：其一，输出功率大于阈值，消耗体力；其二，输出功率等于平衡值，不消耗体力；其三，输出功率低于阈值，恢复体力。这三种状态构成了策略设计的全部素材：赛道上任何一个位置，教练所要做的事情归根到底只是决定让运动员处于哪一种状态、以何种程度处于其中。下面是同时考虑能量供给模型与路况之后的策略选择模型，其基本思想是把赛道上每一处的位置映射为该处应当施加的目标功率，使全程的时间总和最小，同时保证任何时刻的体力状态都不越出可行域。

从形式上说，这是一个以赛道上位置为自变量的最优控制问题：状态变量是 CP、乳酸与氧的浓度，控制变量是输出功率 $P(s)$，约束来自能量供给模型所给出的功率上限与超限机制，目标函数是完赛时间 $\int \mathrm{d}s/v$。由于约束中包含与速度、坡度、曲率相关的非线性项，该问题一般没有闭式解；但正如后文所示，它的最优解在绝大多数路段上都呈现为一个极其简单的形式，即维持阈值速度。

该策略选择逻辑可以概括为一张以体力状态为纵轴、以赛道位置为横轴的示意图：图中把全程划分为若干区间，每一区间标注出运动员所应处于的状态——在下坡与顺风段允许滑行而不输出功率，在直道段以阈值功率稳定推进，在上坡段把功率提高到阈值以上，而在急弯前则短暂地降为零以完成制动。体力曲线随之呈锯齿状起伏，其包络始终不触及代表「耗尽」的下边界，唯有在终点前才被允许短暂地越界以换取冲刺速度。

\subsection{平坦直道上的策略}

我们从最简单的情形出发，即运动员在平坦道路上沿直线骑行。可以想见，维持更高的速度对应着更高的功率输出。此处有两种策略：

\begin{itemize}
\item 运动员以某一特定速度持续运动，在该速度下，输出功率的阈值恰好足以使运动员以恒定速度前进。
\item 运动员在高功率与低功率之间交替，其实际表现是交替地加速与减速。加速过程与减速过程需要保持对称，以使该过程得以成立。
\end{itemize}

对交替加减速的过程，速度变化的加速度可以写为
\[
a=\frac{v_0-v}{m^{*}v_0 v}P_0+\frac{\Delta P}{m^{*}v},
\]
即加速度由两部分构成：一部分来自当前速度偏离参考速度 $v_0$ 所引起的力的不平衡，另一部分来自功率增量 $\Delta P$ 在速度 $v$ 下所提供的额外推进力。第一项在速度高于参考速度时自动变为负值，从而使系统倾向于回到参考速度；第二项则把功率的瞬时增减转换为加速度的正负。两者的叠加刻画了「冲一下、歇一下」这种骑法的动力学本质。

要在两个速度之间交替，需要严格控制时间与加速度等参数，这对运动员而言相对难以实现。另一方面，与维持速度的策略相比，这种交替在一个周期内并不能提供更多的能量：
\begin{equation}
Pt=a\int v^{3}\,\mathrm{d}t+f\int v\,\mathrm{d}t,
\end{equation}
其中左端为总能耗，右端第一项对应空气阻力所消耗的能量、第二项对应与速度成正比的阻力所消耗的能量。

在平均速度保持不变的变速情形下，能量需求反而更高，因此合理的策略是以某一确定的速度持续运动。也就是说，首先以尽可能短的时间加速到某一给定速度，在该速度下运动员恰好能够维持恒定速度。这一结论构成了后续所有策略的基准：在平坦直道上，最优的做法是尽快达到阈值速度并一直保持，而不是不断地冲击与回落。其数学根源在于空气阻力的二次性：若在高于平均值与低于平均值之间摆动，则高速段所增加的阻力功将以平方的速度增长，而低速段所节省的阻力功只能线性地减少，两者相抵之后净损失为正。

\subsection{急弯与丘陵赛道上的策略}

对于弯道，如果速度大于弯道所允许的最大速度，则在进入弯道时减速至该速度并加以维持，而在离开弯道、进入直道时则遵循直道上的策略。也就是说，弯道在策略上表现为一段「速度上限被人为压低」的路段，其前后各需要一个短暂的减速与再加速过程。需要注意的是，减速的时机应当尽量推迟、幅度应当尽量小：因为制动所耗散掉的动能无法回收，而再加速又需要额外的功率，两次代价都会直接转化为时间损失。

从直道进入上坡时有两种方案：

\begin{enumerate}
\item 若上坡路段较长，则可以把上坡视为一段附加了额外阻力的直道。参照上述分析，我们所选择的策略是仅在上坡之前的瞬间减速。
\item 若上坡路段较短，也可以在进入上坡之前的短距离内减速以恢复体力，然后再以尽可能快的速度离开上坡路段。
\end{enumerate}

两种方案的分界取决于上坡的长度与坡度：上坡越长，把其视为「带额外阻力的直道」就越合理，因为此时任何一次冲刺都无法覆盖整段爬坡，反而会透支体力；上坡越短，则越接近一次爆发性用力，此时在坡前稍作恢复、在坡上倾力输出便更为划算。

对于下坡，采用与直道相同的策略。当从下坡进入直道时，采用由直道进入上坡时的策略选择方式。

\section{模型在计时赛赛道上的应用}

在本节中，我们首先对三条赛道进行建模：其中一条是我们自行设计的简单赛道，另外两条是 2021 年的真实赛道。

\subsection{赛道模型}

我们从自行设计的赛道开始，它由直道、坡道、$90^\circ$ 与 $180^\circ$ 弯道构成，其实际形状与各项参数均已标注给出。设计这条赛道的目的是在可控条件下检验模型对各类特征地形的响应，因为真实赛道很难同时具备足够多的急弯与足够陡的爬坡；而一条被刻意设计得「处处刁难」的赛道，恰好可以把策略模型在各个分支上的行为都暴露出来。

这条赛道在平面图上呈近似对称的折返形状：起点之后是一段直道，随后依次经过 $90^\circ$ 弯、上坡、$180^\circ$ 折返弯与下坡，最后沿原路方向返回。其对称性使我们可以直接比较「同一条路段在顺风与逆风两种情形下」的差异，而无需担心地形本身的不对称干扰结论；这一点在后文讨论风的影响时尤为重要。

真实赛道由直线段与圆弧段重建而成，高程变化则通过估计赋值。重建的基本流程如下：先从官方路线图上按等比例提取路线的平面走向；再把路线分解为若干直线段与圆弧段，并逐段记录其长度或半径；随后依据公开的剖面数据为各段赋予坡度；最后把各段首尾相接，形成一条以弧长参数化的闭合或开放曲线。这一流程的误差主要来自两处：一是小半径弯道被简化或被忽略，二是高程数据本身是估计值而非实测值。

考虑到 2021 年奥运会冠军的平均速度约为 13 km/h，我们只统计半径小于 60 m 的弯道，而把半径更大的弯道视为与直道等价。这一处理与上一章关于过弯速度上限的分析相对应：半径足够大的弯道所允许的速度已经高于运动员的实际速度，因而在策略上不再构成约束，把它们并入直道既简化了模型，也不会引入明显误差。

在忽略部分圆弧的条件下，重建模型绕行一圈的长度为 21.1 km，与实际赛道的总长度相比误差约为 4.5\%，这是可以接受的。以同样的方法，我们对在比利时举行的国际自行车联盟世界锦标赛计时赛赛道进行了建模。重建模型的总长度为 43.0 km，与实际赛道 44.3 km 相比误差约为 3\%。

这两条赛道分别对应两种不同的挑战：东京赛道位于城市与赛道混合的路线上，弯道密集、坡度起伏频繁，因而对加减速过程与过弯速度上限极为敏感；比利时赛道更长、更接近典型的公路赛道，直道与缓坡占比较高，因而更考验运动员在长时间内维持阈值功率的能力。把同一位运动员放到这两条赛道上，其最优策略的形态会明显不同，这正是我们选择它们作为检验对象的原因。

两条赛道的重建结果都以「实际路线」与「模型路线」并列的方式呈现：实际路线保留了原始图纸中的全部细节，包括路面的细微弯曲与短小的连接段；模型路线则由长度与半径都被明确记录的直线段与圆弧段拼接而成，去掉了不影响策略的琐碎细节。把二者并置可以看出，重建路线在整体走向、弯道位置与高程趋势上与实际路线一致，只在极短尺度的局部形状上有所简化——这正是误差被控制在 3\%--5\% 之内的原因。

\subsection{模型应用的结果}

四类运动员在三条赛道上的所需时间如下（单位为分钟）：在东京赛道上，男性计时赛专精选手为 28.10、男性冲刺型选手为 31.32、女性计时赛专精选手为 32.69、女性冲刺型选手为 32.72；在比利时赛道上，相应四类运动员分别为 53.02、57.02、54.23 与 58.74；在自行设计赛道上，相应四类运动员分别为 21.01、22.95、22.41 与 23.85。

在东京与比利时的赛道上，我们的模型应用结果略逊于实际世界锦标赛的成绩，但这同时也说明了我们模型的合理性。从结果中可以看出，男性运动员的表现优于女性运动员，计时赛专精选手的表现优于冲刺型选手。以下是这些结果中更为详细的信息。

在这些数字之中，有一处细节值得专门指出：在东京赛道上，女性计时赛专精选手与女性冲刺型选手的成绩分别为 32.69 与 32.72，二者几乎相同；而在比利时赛道上，两者的差距则扩大到 4.51 分钟。这一现象与赛道特性完全吻合——东京赛道弯道密集、起伏频繁，运动员必须反复减速与再加速，而在这类路段上决定成绩的主要是瞬时功率而非续航能力，于是冲刺型选手的短板被大幅削弱；比利时赛道以长直道与缓坡为主，运动员可以长时间保持高速，此时有氧能力便成为决定性的因素。

在我们所建立的模型中，运动员很好地遵循了我们设定的策略。我们的决策以路况为依据，这一点同样反映在功率、速度与赛道上位置的关系之中：在直道段，功率保持在阈值附近而速度基本恒定；进入弯道前功率迅速下降、速度随之降低，出弯后功率短暂抬升以恢复速度；在上坡段功率明显提高而速度下降，在下坡段则相反。当参赛者在该赛道上采用相同的策略时，他们的表现与其基本素质相关。同时我们也可以看到，在地形变化较大的情形下，速度会受到叠加的影响，这也与现实情况相一致。

从这三幅图（分别对应自行设计赛道、东京赛道与比利时赛道）可以看出同一种规律：在纵轴为功率的图中，功率曲线在直道段几乎是一条水平线，在弯道与坡道处则出现尖锐的峰谷；在纵轴为速度的图中，速度的起伏与功率的起伏方向相反，且在地形变化剧烈处被明显放大；而在以赛道上位置为横轴的展开图中，三者的对应关系一目了然——功率的每一次抬升都对应速度的回升，每一次下坠都对应一个弯道或一段陡坡。三幅图的横轴方向即为行进方向，因而可以直观地看出运动员在何处「赚」了时间、在何处「赔」了时间。

\subsection{对目标功率偏离的敏感性分析}

在自行设计的赛道上，我们考察了实际输出功率偏离我们所设定的目标功率时完赛时间的变化。当偏离量为 0 时，完赛时间为 1260.80 s；当偏离量为 $-10\%$ 时，完赛时间为 1272.55 s；当偏离量为 $-20\%$ 时，完赛时间为 1306.60 s；当偏离量为 $-30\%$ 时，完赛时间为 1358.67 s；当偏离量为 $-40\%$ 时，完赛时间为 1445.73 s；当偏离量为 $\pm 20\%$ 时，完赛时间为 1289.00 s；当偏离量为 $\pm 30\%$ 时，完赛时间为 1300.04 s。

由上表可以看出，即使目标功率只有理想目标的 70\%，所造成的时间偏离也仅为 7.76\%；而当目标功率偏离理想目标区间达 30\% 时，所造成的时间偏离也仅为 3.11\%。这些数据表明，我们的策略模型对运动员目标功率的偏离并不十分敏感，这意味着我们的策略在实际使用中具有较好的价值：运动员无需精确到瓦地执行计划，只要把功率控制在阈值附近的一个合理带宽之内，成绩的损失就相当有限。

进一步比较单向偏离与双向偏离两组数据，还可以看出另一层信息：在同样的偏离幅度下，持续偏低（$-20\%$、$-30\%$）所造成的时间损失，明显大于在理想值上下对称摆动（$\pm 20\%$、$\pm 30\%$）所造成的时间损失。这与上一章关于交替加减速的分析完全一致——功率的波动会被空气阻力的二次性放大为能量的净损失，而温和的双向摆动之所以优于单向偏低，是因为高速段所赢回的时间部分抵消了低速段的损失。这一结论对实战的指导意义十分明确：宁可让功率在目标值附近轻微起伏，也不要长时间低于目标值。

\section{天气对骑行的影响}

\subsection{天气的影响}

在自行车运动中，天气与路况都会影响运动员的表现。在各种天气因素中，我们主要关注风、风向与降水状况。

在自行车的直线运动中，风与风向主要影响空气阻力，从而影响运动员在赛道上的实际表现。但在转弯过程中，风向相对于运动员速度的方向在不断变化，因此我们需要讨论风作用下的平均空气阻力。这一区别的实质在于：直道上相对速度是一个恒定的标量 $v\pm v_{wind}$，而在弯道上，运动员的运动方向在持续改变，风速矢量在运动员坐标系中的分解也随之改变，于是相对速度成为转角的函数，必须通过积分才能得到平均阻力。

设自行车需要转过角度 $\gamma$、转弯半径为 $R$，并把风向分解为垂直于运动员初始运动方向与沿着该方向的两个分量，分别记为 $v_x$ 与 $v_y$。

经过时间 $t$ 之后，自行车从原先蓝色方块所示的位置移动到黄色方块所示的位置。在此过程中，功率与阻力的关系、转过的角度以及相对速度分别满足
\begin{equation}
\eta P=p v_r^{2}v,
\end{equation}
\begin{equation}
\varphi=\frac{vt}{r},\qquad \gamma=\frac{vT}{r},
\end{equation}
\begin{equation}
v_r=v+v_x\sin\varphi+v_y\cos\varphi.
\end{equation}
其中第一个式子把功率、相对速度与地速联系起来，第二个式子用转角参数化时间，第三个式子则给出相对速度随转角的变化——垂直于初始方向的风分量以正弦方式进入，沿着初始方向的分量以余弦方式进入，二者在转弯过程中交替地由「侧风」变为「顶风」。

该过程的几何示意可以概括为：蓝色方块表示转弯开始时的位置与航向，黄色方块表示经过时间 $t$ 之后的位置与航向，二者之间的圆弧对应转过的角度 $\varphi$；风速沿两个互相垂直的方向被分解为 $v_x$ 与 $v_y$，随着航向的旋转，这两个分量在运动员坐标系中的投影按正弦与余弦的规律互换，相对速度 $v_r$ 因而随转角连续变化。

于是我们可以通过积分计算平均风阻：
\begin{equation}
\bar{f}_T=p\int_{0}^{T}v_r^{2}\,\mathrm{d}t.
\end{equation}

特别地，我们计算了转弯角度 $\gamma$ 分别为 $\pi/2$ 与 $\pi$ 时的平均风阻，结果为
\begin{equation}
\bar{f}_{\pi/2}=p\left\{v^{2}+\frac{1}{2}\left(v_x^{2}-v_y^{2}\right)+\frac{2}{\pi}\left[2v\left(v_x+v_y\right)+v_xv_y\right]\right\},
\end{equation}
以及
\begin{equation}
\bar{f}_{\pi}=p\left(v^{2}+\frac{4}{\pi}vv_x\right).
\end{equation}
这两个结果揭示了弯道上风的影响结构：对 $90^\circ$ 弯而言，两个风分量都会进入平均阻力，且交叉项 $v_xv_y$ 表明斜向风比同样强度的纯侧风或纯顶风更为不利；对 $180^\circ$ 弯而言，$v_y$ 项在积分过程中被抵消，只剩下与 $v_x$ 成正比的线性修正，也就是说，掉头弯只对某一个方向的风敏感。

在降雨的情形下，水会改变路面的摩擦系数。参照对急转弯的分析可以发现，自行车过弯所允许的最大速度会随之降低，而这将影响运动员的骑行策略。换言之，雨天并不只是让人「感觉更危险」，它在模型中的确切体现是过弯速度上限 $v_m$ 的下降，从而迫使运动员在弯道前更早、更大幅度地减速，并在出弯后付出更多的再加速代价。

\subsection{天气的敏感性分析}

我们以男性专精选手为例，让其在各种风况下完成我们自行设计的赛道，结果如上所述。

根据在风速为 $5\sqrt{2}\ \mathrm{m/s}$ 时绘制的结果可以认为，风向确实会显著影响完成比赛所需的时间。风向对比赛的影响与赛道本身密切相关，这意味着图中所呈现出的对称性是可以解释的，因为我们自行设计的赛道本身具有一定的对称性。若赛道不具备这种对称性，例如某一条赛道由一段长直道与一段密集弯道构成，则风向的影响曲线将呈现出明显的不对称：顺风段的收益无法弥补逆风段的损失，净效果取决于两条路段在总长度中的占比。

在风向固定为 $45^\circ$、风速逐渐增大的情形下，同样可以看出运动员的平均速度受到很大影响；图中括号内的数字分别表示风速与风向。我们知道，风况对车手的影响是一个相对复杂的过程，它不仅与赛场有关，还会影响车手的策略选择。其复杂性体现在两个层面：其一，风改变了同一功率所对应的地速，从而改变了各路段之间的时间分配；其二，风还改变了阻力的空间分布，使得原本均匀的最优功率分配出现偏向——在逆风段略微降低功率、在顺风段略微提高功率，往往能比全程恒功率取得更好的总时间。

风况敏感性分析的结果以两条曲线的形式呈现：一条是完赛时间随风向变化的曲线，在 $0^\circ$ 至 $360^\circ$ 的范围内呈周期性的双峰双谷，谷底对应逆风最少的航向，峰值对应逆风最多的航向；另一条是完赛时间随风速变化的曲线，随着风速由零逐渐增大，完赛时间单调上升，且上升的斜率随风速增大而变陡——这与空气阻力对相对速度的二次依赖完全一致。图中括号内标注的数值分别给出该数据点所对应的风速与风向，以便逐点核对。
\section{团体计时赛的修正模型}

\subsection*{附加假设}

要把我们的个人计时赛模型推广到团体计时赛，需要补充更多假设。这是因为团体计时赛引入了两个新的机制：一是队员之间的气动遮蔽，它使不同位置的运动员面对不同的阻力；二是队员之间的体力调度，它使个人的体力可以在队伍内部转移。

\begin{itemize}
\item 车队全部六名成员具有相同或相近的身体素质。若车手之间的身体素质存在差异，问题将变得复杂得多。在这一假设之下，队伍可以被视为六个能力相同的个体所组成的整体，其策略问题便归结为「何时由谁领骑」的调度问题，而不必再处理「强手是否应当多承担」这一额外的分配问题。
\item 在队列中，队员排列紧凑。我们认为具体的交接过程近似等价于忽略若干因素之后的结果，这些因素包括从前部退到队尾所花费的时间，以及由此产生的加速、减速与位置变化。换言之，我们把轮转视为瞬时完成的身份交换，从而可以专注于领骑时长的分配，而不必刻画交接瞬间的动力学细节。
\end{itemize}

\subsection*{轮转领骑}

在团体计时赛中，少数运动员可以在车队前方领骑，而其他队友则在其后的尾流中跟骑。领骑的运动员随后轮转到队尾。跟骑效应最多可以把整体空气阻力降低至原来的 $2/3$[15]。这一效应的物理实质是：领骑者在前方排开空气，在其后方形成一个压力较低、湍流较强的尾流区，处于其中的运动员所面对的相对风速显著下降，于是同一速度所需的功率随之减少。

我们假设，当一名运动员在领骑时，其后方运动员的空气阻力降低为原来的 75\%。当领骑者以高于其阈值的输出功率前进时，其后五名车手能够以低于其阈值的输出功率恢复体力。这一组假设把团体计时赛的核心矛盾刻画得十分清楚：队伍的整体速度由领骑者决定，而队伍的持续能力由跟骑者的恢复速度决定；前者要求领骑者尽力，后者要求领骑者不要过久地留在前方。

\subsection{冲刺阶段之前的团队策略}

其策略是在前方轮流骑行一定的距离，直到体力接近耗尽，然后退到队尾恢复，从而使整个队伍实际上始终比单个车手更快。设车队的整体行进速度为 $V$，则根据运动模型，可以得到第一名车手与其他车手的功率需求分别为 $P_0$ 与 $P_1$：
\begin{equation}
\eta P_0=V\left(f_{air}+f_{other}\right),
\end{equation}
\begin{equation}
\eta P_1=V\left(0.75f_{air}+f_{other}\right),
\end{equation}
其中 $f_{air}$ 为空气阻力，$f_{other}$ 为其余各类阻力之和。两式之差即跟骑所带来的功率节省，它正比于 $0.25f_{air}V$，因此速度越高、空气阻力占比越大，跟骑的收益就越显著——这也解释了为何团体计时赛的收益远大于爬坡赛段。

根据能量供给模型，可以确定车手在 $P_0$ 输出功率的条件下能够坚持时间 $T_0$，而完全疲劳的车手在 $P_1$ 输出功率的条件下完全恢复需要时间 $T_1$。为实施轮转领骑的策略，必须满足 $T_1\leqslant 5T_0$。这一不等式的含义是：一名领骑者在前方耗尽其体力所需的耐久时间，必须不短于队尾五名队员依次完成恢复所需的时间；只有在这样的节奏下，队伍才不会因为某位队员尚未恢复而被迫降速。在这一限制之下，我们可以求出能够满足要求的最大速度。

对于弯道、上坡与下坡路段，我们认为车队在这些路段上的领骑过程与在直道上相同。把车队视为一个整体，其前进策略与个人在这些路段上的策略一致。这一处理的合理性在于：队形在这些路段上虽然会拉长或压缩，但遮蔽关系依然存在，阻力结构的相对比例不变，因而策略的形态也保持不变。

\subsection{最终冲刺阶段的团队策略}

我们提出一种放弃车手体力恢复的策略：提高车队的整体速度，使第一名车手之后的每位车手都无法恢复体力，即 $P_1$ 等于阈值输出功率。这一策略的逻辑是：既然终点已经不远，此后不再需要「可持续」，只需要「尽快」，那么保留恢复能力便失去了意义，把全部剩余体力转换为速度才是最优的。

在第一名车手耗尽体力、甚至超出其功率输出曲线之后，他将放弃追赶车队。第二名车手随后重复第一名车手的策略，并从车队中突围而出。如此，队伍实际上被拆解为一连串的「自杀式领骑」：每一位领骑者都在自己还能承受的极限内把速度推到最高，然后退出，把位置让给下一位体力尚存的队友。

这一冲刺阶段的队形变化可以概括为一条「逐级断裂」的示意：最初六人成一条直线，第一名车手居于最前，其后五人的体力条依次缩短；随着时间推移，最前方的一格被耗尽并脱离队形，队列整体前移一位，原本排在第二位的车手接管领骑；如此重复四次，最终只剩一名体力最充沛的车手以最高速度冲过终点。图中每一格的位置与该车手当时的体力水平相对应，断裂的顺序即领骑的轮转顺序。

设车手与终点线之间的距离为 $L$，最后四名选手消耗各自的体力以进一步加速至速度 $v$。设第一名车手的功率水平使其体力以速率 $U_0$ 消耗，而后三名车手的相应速率为 $U_1$；他们剩余的体力分别为 $3V$、$2V$、$V$ 与 $0$，而超出能量供给曲线的体力消耗上限为 $M$。设四名车手各自的领骑时间分别为 $t_i\ (i=1,2,3,4)$，车队在六人直线阶段的阈值速度为 $v_0$，则可以得到如下公式：
\begin{equation}
\begin{cases}
V(i-1)+M=u_1\displaystyle\sum_{j\neq i}t_j+u_0t_i,\\[1em]
v\displaystyle\sum_{i=1}^{4}t_i=L.
\end{cases}
\end{equation}
其中第一组方程的含义是：对第 $i$ 名领骑者而言，他在登场时所拥有的体力（基础存量 $V(i-1)$ 加上允许透支的上限 $M$）恰好等于他在其余队员领骑期间以速率 $u_1$ 恢复的总量，加上他自己领骑期间以速率 $u_0$ 消耗的总量；第二组方程则只是说，四段领骑所覆盖的距离之和等于剩余距离 $L$。

通过调整决策变量 $v$ 与 $t_i\ (i=1,2,3,4)$ 的取值，使
\[
\frac{L}{v_0}-\sum_{i=1}^{4}t_i
\]
取得最大值，即表示四人冲刺阶段达到最佳效果。该目标函数的含义十分直观：第一项是若全队以阈值速度 $v_0$ 通过距离 $L$ 所需的时间，第二项是采用冲刺策略实际花费的时间，二者之差即为团队策略在该路段上所赢得的时间收益。当这一差值取到最大时，冲刺阶段的队形安排、每位领骑者的时长与整体的行进速度便同时被确定下来。

\subsection{团体计时赛模型的求解}

利用我们在个人计时赛中所设定的参数，可以根据上述策略求解团体计时赛中的提升幅度。结果表明，在正常推进阶段，车队的平均速度相比个人提高了 14.43\%；在冲刺阶段，平均速度相比个人提高了 18.52\%。

这两个数字之间的差别同样值得注意：冲刺阶段的提升幅度大于推进阶段，其原因在于冲刺阶段允许透支体力，因而跟骑所带来的功率节省可以被更彻底地转化为速度；而推进阶段必须维持可持续的节奏，提升幅度因此受到恢复能力的约束。这一结果与团体计时赛的实际经验相符——团队的最大优势往往出现在最后的十公里。

\section{模型评价与进一步讨论}

\subsection{模型评价}

\subsubsection{模型的优点}

\begin{itemize}
\item 问题被划分为不同的模块，并综合运用了物理学与生物学的相关知识。能量供给、运动学与策略优化三部分彼此独立又相互衔接，任一模块都可以单独替换或细化，而不必重写整个模型。
\item 从基础的能量供给模型出发，我们的建模更贴近现实。我们没有直接对功率曲线作经验拟合，而是从物质守恒与生理限制出发推导出功率曲线，因此模型能够解释而不仅是复现观测到的现象。
\item 在不同运动员之间进行了全面的比较，模型结果具有说服力。四类运动员、三条赛道的组合覆盖了性别与能力结构两个维度，结果与真实比赛成绩的一致性为模型提供了外部验证。
\end{itemize}

\subsubsection{模型的缺点}

\begin{itemize}
\item 在模型建立过程中，许多内容被简化，这可能使模型对实际问题的反映变得不够准确。例如把运动员视为质点、把弯道内的运动视为匀速圆周运动、把传动效率取为常数，都会在不同程度上抹去真实骑行中的细节。
\item 在策略选择方面，分析尚不充分，讨论只说明了所选择的策略是较优策略。我们并未严格证明该策略在全局意义下最优，因而严格说来结论应被理解为「一种表现良好的可行策略」，而非「唯一的最优解」。
\end{itemize}

\subsection{模型可能的改进}

在模型求解时曾有使用变分法的设想，但由于时间限制未作进一步尝试。事实上，有可能利用变分法把功率作为距离的函数来求解，这对于精确描述该问题将很有帮助。若把完赛时间写为
\[
T=\int_{0}^{S}\frac{\mathrm{d}s}{v\left(s\right)},\qquad v=v\left(P\left(s\right),s\right),
\]
并附加以体力状态为约束的微分方程，则该问题具有古典变分问题的形式，其欧拉—拉格朗日方程可以直接给出最优功率随位置变化的解析结构。这将是后续工作中最值得推进的方向。

\section*{参考文献}
\begin{enumerate}
\item History of Road Cycling. [Online]. Available: https://olympics.com/en/sports/cycling-road/\#discipline-history-of
\item D. Lange, ``Participants in road/paved surface bicycling in the U.S. 2011-2020,'' Statista, Tech. Rep., Jun. 2021. [Online]. Available: https://www.statista.com/statistics/763746/road-paved-surface-bicycling-participants-us/
\item CGavin. Fuji Speedway (FSW), the cycling course of Tokyo 2020. [Online]. Available: https://720yun.com/t/267jvzwwzu3?scene\_id=20194410
\item Wikipedia. Team time trail. [Online]. Available: https://en.wikipedia.org/wiki/Team\_time\_trial
\item [Online]. Available: https://zhuanlan.zhihu.com/p/37938089
\item Damian. Power profile in cycling. [Online]. Available: https://cyklopedia.cc/cycling-tips/power-profile/
\item J. Fortson. Functional threshold power: What ftp means to cyclists. [Online]. Available: https://www.trainerroad.com/blog/what-ftp-really-means-to-cyclists/
\item Damian. Vo2 max training, ranges, workouts. [Online]. Available: https://cyklopedia.cc/cycling-tips/vo2-max/
\item J. Fortson. How to use your cycling power curve to find your strengths and weaknesses. [Online]. Available: https://www.trainerroad.com/blog/how-to-use-your-power-curve-to-find-your-strengths-and-weaknesses/
\item J. Lighthill, \emph{An informal introduction to theoretical fluid mechanics}, Jan. 1986, pp. 97--121. [Online]. Available: https://www.osti.gov/biblio/5126674
\item A. Denham. (2019, Jun.) Drivetrain Efficiency: What's The Difference In Speed Between 1X and 2X? [Online]. Available: https://www.cyclingabout.com/drivetrain-efficiency-difference-speed-between-1x-2x/
\item N. R. C. U. T. R. B. C. for the National Tire Efficiency Study, \emph{Tires and passenger vehicle fuel economy: informing consumers, improving performance}. Transportation Research Board, 2006, vol. 286, p. xvi.
\item G. de Groot, A. Sargeant, and J. Geysel, ``Air friction and rolling resistance during cycling,'' \emph{Medicine and science in sports and exercise}, vol. 27, no. 7, pp. 1090--1095, Jul. 1995.
\item A. Pytel and J. Kiusalaas, \emph{Engineering Mechanics: Dynamics - SI Version}. Cengage Learning, 2010, p. 654. [Online]. Available: https://books.google.co.jp/books?id=8gYKzgEACAAJ
\item Wikipedia. Drafting (aerodynamics). [Online]. Available: https://en.wikipedia.org/wiki/Drafting\_(aerodynamics)
\end{enumerate}

\section*{给车手的比赛指导}

在计时赛赛道中，运动员按固定间隔分别出发，并以其到达终点的时间排名。规则很简单，但要取得理想的成绩却很难。它受到诸多因素的影响，包括运动员的身体素质、对赛道的了解、天气以及比赛中采用的策略。本文正是为你——车队体育主管——而写的，以期为你在帮助车手取得更好成绩方面提供指导。

\subsection*{第一，充分掌握车手的身体素质}

请关注车手日常训练的数据，它将描摹出他们的功率曲线，由此你可以判断他是爬坡手、冲刺手还是其他类型。你还需要经常查阅这张图表，以观察运动员身体能力的变化趋势。

在功率曲线中，第一部分显示车手的爆发力，在这一部分水平较高的车手往往专精于冲刺。我们把在第二部分表现优异的车手称为爬坡手，这意味着他们专精于包含多段长爬坡的比赛。而第三部分水平较高则意味着车手拥有出色的耐力，他们在长距离比赛中表现优异。

功率曲线由 ATP-CP、糖解与有氧三个部分混合而成，它们共同构成了车手的供能系统。通过分析该系统中各类物质的变化，我们可以详细了解车手的身体状态，并在此基础上为你提供专业建议。

\subsection*{第二，提前了解前方路况}

比赛开始之后，如果车手对前方的道路与天气一无所知，就很难作出正确的决策。车手的速度、加速度与其功率输出之间的关系，一般由环境所决定。

以下是你应当了解的、关于道路与天气如何影响车手表现的要点：

\begin{itemize}
\item 在上坡路段，车手需要更大的功率来维持速度；在下坡路段，所需的功率较小。
\item 当车手通过弯道时，请记住绝不要维持过高的速度，尤其是在下雨时。
\item 你还应当明白，相对于风的速度越快，你所受到的空气阻力就越大。
\end{itemize}

\subsection*{第三，选择最佳策略}

好的策略往往能把车手引向成功。基于我们模型的结果，我们在此给出关于车手在比赛中应当采用何种策略的建议。

\begin{itemize}
\item 在道路上，最重要的事情是明确阈值速度。但如何找到它呢？只需不断加速，直到车手达到这样一个最大速度：它既不会让他感到十分疲惫，也不会让他摔倒在弯角处。这比交替进行加速与减速的过程要好得多。我们的模型也解释了其中的原因：功率的波动会被空气阻力的二次性放大为额外的能量消耗，因而恒速骑行在绝大多数路段上都是最优的。
\item 当车手接近转弯处时，首先评估半径，以估计不发生侧滑的速度。用刹车迅速减速，车手便可以在转弯过程中得到短暂的休息。
\item 在长坡路段，基本策略与平坦道路上相同。但要注意，当车手从低阻力区域移动到高阻力区域时，例如从下坡区域移动到上坡区域，要评估车手能否冲过这一区域。如果你对车手的体力有把握，就让他冲过去；否则就告诉他减速。
\item 在加速时，车手需要估计到目的地的距离，并确保自己不会在体力耗尽之前无法抵达。
\item 在团体计时赛中，轮转领骑非常重要。让你的车手轮转领骑，这将使他们表现得比各自为战时更好。
\end{itemize}

\end{document}
